% =============================================================================
% methods.tex
% Rewritten "Adaptive Gating Policy" section.
% Assumes Section~\ref{sec:envs} contains the SMDP formalisation, so this
% section focuses on the gate itself: motivation, observations, architecture,
% training.
%
% Notes:
%   - References eq:smdp from intro.tex / Section 3.1.
%   - References Figure~\ref{fig:scaling} (intro / appendix).
% =============================================================================

% =============================================================================
\section{Adaptive Gating Policy}
\label{sec:method}
% =============================================================================

The gating policy is a lightweight network trained on top of a frozen
AlphaZero planner that selects, at each meta-step, a planning budget
$K \in \Kset$.
We described the variable-delay setting it operates in formally in
Section~\ref{sec:envs}; here we describe the policy itself---what it
observes, how it is structured, and how it is trained.

The empirical scaling result in Figure~\ref{fig:scaling} characterises
the tradeoff the gating policy must manage.
Across all five environments, additional MCTS simulations produce
monotonically better actions while also taking proportionally longer to
run, with the relationship approximately linear over the simulation
ranges we study.
This co-scaling makes the allocation problem well-posed but non-trivial.
It is well-posed because each unit of compute buys a roughly predictable
unit of quality.
It is non-trivial because the value of that quality varies sharply with
state: a position where most sensible actions lead to similar outcomes
derives little benefit from refined search, while a position where the
right action is several rollouts deep can swing the whole episode.
The natural response is to allocate compute where the marginal quality
gain is high and skip it where it is not.
We do not attempt to learn this tradeoff jointly with the planner.
The planner already represents the action-quality side of the tradeoff
through its visit-count distributions; our task is only to learn, on top
of it, a state-dependent decision about how long to run it.
This separation has two practical advantages: it reuses a strong base
planner without compromising its training, and it isolates the
meta-reasoning problem in a small network whose forward pass is
microseconds rather than milliseconds.

\subsection{Observations and architecture}

At each meta-step the gating policy takes three inputs: (1) the raw
game grid, (2) the planner's intermediate spatial features, extracted
from its frozen trunk and used as a compact summary of what the planner
currently infers about the position, and (3) the planner's scalar value
estimate $V(s_t)$.
In clock environments the observation additionally includes the
remaining time budget for both players.
Together these inputs give the gating policy a direct view of the
board, a compressed summary of the planner's confidence, and---where
applicable---the time pressure conditioning the decision.

The trunk features matter because the planner has already done
substantial spatial reasoning by the time the gate is called: extracting
that reasoning costs nothing extra, since it is a side product of the
same forward pass that produces $V$, and it gives the gate a richer
signal than the raw grid alone.
We deliberately do not feed the planner's policy logits to the gate.
Doing so would essentially leak the planner's preferred action into the
budget decision, and we want the gate to learn from structural features
of the state---threat geometry, board density, time pressure---rather
than from a shortcut summary of what the planner already wants to do.

A lightweight network processes these inputs and produces a categorical
distribution over $K$ together with a scalar baseline value used during
training.
Architecture details, including the recurrent variant used for Speed
Hex, are in Appendix~\ref{app:implementation}.

\subsection{Training}

Training proceeds in two phases.
We first train AlphaZero-style base planners for each environment using
standard expert iteration (Appendix~\ref{app:az_primer}) and select
frozen base checkpoints by cross-evaluation across delay budgets
$K \in \Kset$ (Appendix~\ref{app:cross_eval}).
We then freeze the selected planner and train the gating policy with
PPO~\citep{schulman2017proximal} on top of it.

Rollouts are collected over full episodes.
At each meta-step the gate samples $K_t \sim \pigate(\cdot \mid s_t)$,
the agent executes the resulting option---$K_t{-}1$ filler actions
followed by the planner's chosen action---and the meta-reward $R_t$
defined in Section~\ref{sec:envs} is accumulated.
Variable-duration advantages with per-meta-step discount $\gamma^{K_t}$
are computed before each round of PPO updates, adapting Generalized
Advantage Estimation~\citep{schulman2015high} to the SMDP introduced in
equation~(\ref{eq:smdp}).
Only the gating network is updated; no gradient flows to the planner.

Making this meta-MCTS-AlphaZero stack computationally feasible requires
careful implementation: in particular, evaluating all $|\Kset|$ option
transitions in parallel at each meta-step preserves static shapes and
allows the entire training loop to compile to a single XLA program.
Without this trick, the cost of nesting MCTS inside an outer RL loop
dominates wall-clock time.
The full JAX implementation---including the use of \texttt{lax.scan}
for both MCTS inner loops and PPO rollouts, \texttt{vmap} for
environment parallelism, and \texttt{pmap} for device parallelism---is
described in Appendix~\ref{app:implementation}.
